\section{来源审计与核心问题：稀疏性究竟扩展了什么}

本讲的标题很容易诱导出一个错误印象：仿佛只要把模型参数写成“万亿”，就等于让每个 token 都经过万亿参数计算。讲座真正讨论的是另一件事：让模型拥有很大的\textbf{可选参数池}，但每个 token 只激活其中少量专家，从而把总参数量与单 token 计算量部分解耦。理解这一点，是读懂后续所有速度、质量和硬件图表的前提。

本次重写使用 Stanford Online 官方录像、官方手动英文字幕、38 张从 1080p 录像中恢复并人工筛选的完整教学页，以及 Switch Transformer、GShard、T5、mT5 与 V-MoE 原始论文。录像没有公开独立 slide PDF，因此每张图都标注视频时间；重复动画、片头和 Q\&A 中反复跳回的同一页不重复插入，但口头解释保留在正文中。

\lecturefigure{slide-01-title.jpg}{讲座主题：Scaling Transformers through Sparsity。}{Stanford Online 官方视频 00:00:06--00:01:41。}

这张标题页之后，讲者没有直接进入路由公式，而是先回到 dense scaling 的经验事实：语言模型的 loss 随参数、数据和计算预算呈现可预测的幂律趋势。MoE 的研究问题并不是否认 dense scaling，而是问：能否在保持每个样本计算近似不变的情况下，再增加一个“条件使用参数”的扩展维度？

\lecturefigure{slide-02-neural-scaling-laws.jpg}{Dense neural scaling laws：更大的模型具有更高样本效率，并存在随预算平滑变化的最优规模。}{Stanford Online 官方视频 00:01:42--00:02:35。}

可以用一个简化式表达 dense scaling 的经验形状：

\[
\mathcal{L}(N)=\mathcal{L}_{\infty}+aN^{-\alpha},\qquad a>0,\ \alpha>0.
\]

其中，$\mathcal{L}(N)$ 表示参数规模为 $N$ 时的验证损失；$\mathcal{L}_{\infty}$ 是在该数据分布与建模假设下的不可约下界；$a$ 控制纵向尺度；$\alpha$ 决定扩大模型带来的边际收益衰减速度。这个式子描述的是经验趋势，不是保证任意架构、数据或硬件上都沿同一曲线外推。

\teachervoice{老师先强调 scaling law 的意义：它给研究者一种“扩大规模仍可能获得收益”的信心，但原论文主要研究 dense model。于是本讲把 sparsity 作为新的结构轴，而不是把 MoE 讲成与 scaling 无关的局部模块。}

\lecturefigure{slide-03-new-axis-sparsity.jpg}{新的扩展轴：增加稀疏使用的参数，并尝试与每个样本的计算量解耦。}{Stanford Online 官方视频 00:02:36--00:03:05。}

\begin{importantbox}{资源核算的中心不变量}
讨论 sparse model 时必须同时报告至少五个量：\textbf{总参数量}、\textbf{每 token 激活参数量}、\textbf{每 token FLOPs}、\textbf{跨设备通信量}与\textbf{wall-clock 时间}。只说“参数更多但计算不变”是不完整的，因为未激活参数仍占设备或主存，路由还会产生 all-to-all 通信与负载不均。
\end{importantbox}

\begin{warningbox}{万亿参数不等于万亿参数都被执行}
Switch-C 的 1.6T 是可用参数池的总规模，不是每个 token 的 active path。反过来，active FLOPs 较小也不代表系统免费：存储、通信、容量缓冲、编译形状和低吞吐推理都会暴露额外成本。
\end{warningbox}

\begin{table}[H]
\centering
\small
\caption{读 sparse scaling 时必须分开的资源量。}
\begin{tabularx}{\textwidth}{L{0.22\textwidth}X X}
\toprule
量 & 它回答的问题 & 常见误读 \\
\midrule
总参数量 & 模型一共存储多少可学习权重？ & 把它当作每 token 计算量。 \\
激活参数量 & 一个 token 实际经过哪些 expert 权重？ & 忽略不同 token 的路径不同。 \\
FLOPs & 理论算术工作量是多少？ & 用 FLOPs 直接替代真实延迟。 \\
通信量 & token dispatch 与结果汇聚传输多少数据？ & 忽略 all-to-all 成为瓶颈。 \\
wall-clock & 在指定硬件、batch 与实现下多久达到目标质量？ & 脱离硬件和吞吐场景泛化。 \\
\bottomrule
\end{tabularx}
\end{table}

\subsection{本章小结}

本讲扩展的是\textbf{条件可用参数}，不是无条件增加每个 token 的计算。后文所有结果都应在总参数、激活参数、FLOPs、通信和时间这五个坐标中解释。

